\section{Probabilidades}

\subsection{Fundamentos de Probabilidade}
Ao observar, registrar ou extrair uma amostra de dados a partir de um fenômeno aleatório,
isto é, cujo resultado não pode ser deterministicamente previsto, a Estatística, dentre
outras coisas, procura entender a frequência de ocorrência dos valores amostrados,
que nos dão diversas opções de como analisar seu comportamento. Essas frequências são
estimativas das \textbf{probabilidades} dos eventos observados. A partir de hipóteses
adequadas, podemos modelar as frequências por modelos probabilísticos que, por sua vez,
são objetos de estudo da Teoria das Probabilidades.

A Teoria de Probabilidades encontra seus fundamentos em alguns campos da matemática pura,
tais como Matemática Discreta, Análise e Teoria da Medida. No entanto, neste texto,
estaremos mais interessados na visão da teoria para aplicações estatísticas, como a já
citada modelagem de frequências e estimação de parâmetros.

A noção primitiva que iremos usar como base será de experimento ou fenômeno aleatório,
que modela qualquer fenômeno cujos resultados são conhecidos, porém que podem ser
distintos apesar da repetição de condições iniciais. Com isso, daremos a definição do
espaço amostral

\subsection{Axiomas de Probabilidade}
\begin{def:espaco amostral}
O espaço amostral $\Omega$ de um experimento aleatório é o conjunto de todos os resultados
possíveis do experimento. Um subconjunto de $\Omega$ é chamado de evento.
\end{def:espaco amostral}
Com isso, temos a definição de probabilidade de um resultado, ou subcojunto de resultados.
\begin{def:probabilidade}
A probabilidade $P(E)$ de um evento $E$ num espaço amostral $\Omega$ satisfaz:
\begin{enumerate}
	\item $P(E) \in [0, 1]$

	\item $P(\Omega) = 1$

	\item Se $E_1, \ldots, E_n$ são eventos disjuntos em $\Omega$, também ditos mutualmente
	      exclusivos, então
	      \[
		      P\left(\bigcup_{i=1}^n E_i\right) = \sum_{i=1}^n P(E_i)
	      \]
\end{enumerate}
\end{def:probabilidade}
Probabilidades são frequentemente mostradas como porcentagens e são entendidas como as
chances de um resultado, ou cojunto deles, ser observado num experimento aleatório.

\subsection{Eventos Equiprováveis}
Numa noção intuitiva de que quando o evento não possui vícios, os seus resultados
são ditos como equiprováveis, como os lados de um dado ou as faces de uma moeda honestos.

\begin{def:evento equiprovavel}
Um evento $E \subset \Omega$ é equiprovável se, e somente se, todo subconjunto unitário de
$E$ possui a mesma probabilidade.
\end{def:evento equiprovavel}

\begin{prop:equiprovavel}
A probabilidade de um evento $E$ num espaço $\Omega$ equiprovável
é igual a
\[
	\frac{|E|}{|\Omega|}
\]
\begin{proof}
	Com efeito, temos que
	\[
		P(E_1) = \ldots = P(E_n)
	\]
	em que $E_i$ são eventos unitários de $\Omega$, tal que $|\Omega| = n$. Como eles são
	disjuntos, então a soma de suas probabilidades deve ser igual a 1, ou seja,
	\[
		\forall i \in [n],\ \left(P(E_i) = \frac{1}{n}\right).
	\]
	Segue que
	\[
		P(E) = |E| \cdot \frac{1}{n} = \frac{|E|}{|\Omega|}
	\]
\end{proof}
\end{prop:equiprovavel}

\begin{cor:pie}
Se $\mathcal{F} = \{E_1, \ldots, E_m\}$ é uma família de eventos quaisquer dum espaço
$\Omega$ equiprovável, então
\[
	P\left(\bigcup_{E \in \mathcal{F}}E\right) = \sum_{k=1}^m (-1)^{k+1}
	\sum_{D \in \Parts{\mathcal{F}}\land |D| = k} P\left(\bigcap_{E \in D} E\right)
\]
\begin{proof}
	A cardinalidade da união dos conjuntos de $\mathcal{F}$ é dado pelo Princípio da
	Exclusão e Inclusão, ou seja
	\[
		\left|\bigcup_{E \in \mathcal{F}}E\right| = \sum_{k=1}^m (-1)^{k+1}
		\sum_{D \in \Parts{\mathcal{F}} \land |D| = k} \left|\bigcap_{E \in D} E\right|
	\]
	Dividindo a expressão acima por $n = |\Omega|$ obtém-se a probabilidade da união dos
	conjuntos de $\mathcal{F}$. Contudo, uma vez que a interseção de eventos é um evento,
	então podemos obter a probabilidade da união como
	\[
		P\left(\bigcup_{E \in \mathcal{F}}E\right) = \sum_{k=1}^m (-1)^{k+1}
		\sum_{D \in \Parts{\mathcal{F}}\land |D| = k} \frac{1}{n}\left|\bigcap_{E \in D} E\right|
	\]
	que equivale a expressão desejada.
\end{proof}
\end{cor:pie}

\subsection{Probabilidade como Função Contínua em um Conjunto}

\begin{def:eventos encaixados}
Uma sequência ${E_i}_{i \in \mathbb{N}}$ é encaixado se
\[
	E_1 \subset E_2 \subset \ldots
\]
ou
\[
	E_1 \supset E_2 \supset \ldots.
\]
O primeiro caso é chamado de crescente, enquanto o segundo de descrescente.
\end{def:eventos encaixados}

\begin{prop:eventos encaixados}
Se ${E_i}_{i \in \mathbb{N}}$ eventos encaixados crescentes, então
\[
	\lim_{n \to \infty} P(E_n) = P\left(\lim_{n \to \infty} E_n\right)
\]
\begin{proof}
	.
\end{proof}
\end{prop:eventos encaixados}
